\clearpage
\chapter*{ABSTRACT}
\phantomsection
\addcontentsline{toc}{chapter}{ABSTRACT}
\setstretch{1.5}

\textit{The advancement of Internet of Things (IoT) technology has enabled the
development of sensor-based monitoring systems continuously connected to
cloud platforms. This final project develops a smart IoT-based temperature
monitoring system using ESP32-S3 and DS18B20 sensor simulated on the Wokwi
platform. Temperature data is sent to ThingSpeak via a virtual Wi-Fi
connection every 15 seconds. The system is equipped with an artificial
intelligence module based on Random Forest Regressor running in Python. The
AI model reads historical data from ThingSpeak via a public channel,
constructs seven time-series features, and predicts the temperature 1 minute
ahead (4 steps forward).
The predicted temperature is classified into three categories: cold
($<$25\textdegree C), normal (25\textdegree C -- 30\textdegree C), and hot
($>$30\textdegree C). The AI output is sent back to ThingSpeak and displayed
on the dashboard alongside the actual temperature. Model evaluation uses MAE,
RMSE, and R$^2$ Score metrics. Test results show that the system successfully
performs sensor data acquisition, end-to-end IoT communication, AI-based
temperature prediction, and automatic dashboard visualization.}

\vspace{0.7cm}
\noindent\textbf{\textit{Keywords:}} \textit{smart IoT, ESP32-S3, DS18B20,
ThingSpeak, Random Forest Regressor, temperature prediction, Wokwi.}
